\documentclass[11pt]{article}

% --- Кодировка и язык (pdflatex) ---
\usepackage[T2A]{fontenc}
\usepackage[utf8]{inputenc}
\usepackage[english,russian]{babel}

% --- Математика и оформление ---
\usepackage{amsmath,amssymb,amsthm}
\usepackage{bm}
\usepackage[margin=25mm]{geometry}
\usepackage{enumitem}

\newcommand{\vth}{\bm{\vartheta}}
\newcommand{\Deff}{D_{\mathrm{eff}}}
\newcommand{\Dphys}{D_{\mathrm{phys}}}
\DeclareMathOperator*{\argmin}{arg\,min}
\DeclareMathOperator{\AIC}{AIC}

\title{\vspace{-1.5em}Постановка задачи: подбор модели и оптимизация параметров\\
\large Извлечение параметров проволочных поляризаторов из данных THz-TDS}
\author{Проект THz-Unified-Optimizer}
\date{\today}

\begin{document}
\maketitle
\vspace{-1.5em}

\noindent
Документ фиксирует \emph{математическую} постановку двух связанных задач проекта:
(i)~\textbf{оптимизации параметров} выбранной прямой модели и (ii)~\textbf{подбора модели}
(структурного выбора) среди конкурирующих физических гипотез. Обозначения едины для всех
последующих разделов статьи.

\section{Данные и наблюдаемая величина}
Характеризация проволочного поляризатора (WGP) сводится к измерению комплексных спектров
пропускания методом THz-TDS. Наблюдения — комплексные амплитуды поля
\[
  y_{ij} \;=\; E_{\mathrm{obs}}(\theta_j,\nu_i)\in\mathbb{C},
  \qquad i=1,\dots,N_\nu,\quad j=1,\dots,N_\theta,
\]
на сетке частот $\{\nu_i\}$ и углов ориентации оси пропускания $\{\theta_j\}$. Введём сквозной
индекс $k=(i,j)$, $k=1,\dots,N$, $N=N_\nu N_\theta$, и известную (оценённую по фону/повторам)
дисперсию шума $\sigma_k^2$ с весами $w_k=1/\sigma_k^2$.

\section{Семейство прямых моделей}
Прямая модель — отображение параметров в предсказание наблюдаемой,
\[
  \hat y_k \;=\; f\!\left(\nu_i,\theta_j;\,\vth,\,M\right),
  \qquad \vth\in\Theta_M\subset\mathbb{R}^{p_M},
\]
где $M\in\mathcal{M}$ — \emph{структура} модели (набор включённых физических каналов), а
$\vth$ — её $p_M$-мерный вектор параметров. Физическое ядро — аналитическая модель Бланко
(эквивалентные схемы), дающая частотно-зависимые коэффициенты $t_\perp(\nu;\vth)$ и
$t_\parallel(\nu;\vth)$ в формализме матриц Джонса; в простейшем «один реальный WGP\,+\,идеальный
анализатор»
\[
  \hat E(\theta,\nu) \;=\; \cos^2\!\theta\;t_\perp(\nu;\vth)\;+\;\sin^2\!\theta\;t_\parallel(\nu;\vth).
\]
Кандидатные структуры $M$ различаются включением поправок: проводимость Друде, степенное
рассеяние $\nu^{\gamma}$, эффективный диаметр проволоки $\Deff$, а также явная утечка
параллельного канала
\[
  t_\parallel \;\to\; t_\parallel + \eta(\nu)\,e^{\,j(\delta_0+2\pi\nu\,\tau_{\mathrm{leak}})}.
\]

\section{Оптимизация параметров}
Комплексная невязка и целевая функция (взвешенный нелинейный МНК, эквивалентный
гауссовскому максимуму правдоподобия):
\[
  r_k(\vth)\;=\;y_k-\hat y_k(\vth),
  \qquad
  \chi^2_M(\vth)\;=\;\sum_{k=1}^{N} w_k\,\bigl|r_k(\vth)\bigr|^2 .
\]
Оценка параметров при фиксированной структуре $M$:
\[
  \boxed{\;\hat{\vth}_M \;=\; \argmin_{\vth\in\Theta_M}\;\chi^2_M(\vth)\;}
\]
Задача невыпуклая; итерации Гаусса--Ньютона / Левенберга--Марквардта сходятся к локальному
минимуму, поэтому для мультимодальных (в частности угловых) параметров \emph{обязателен
мультистарт} — иначе значимость завышается. При гауссовском шуме
$\ln L_M(\vth)=-\tfrac12\chi^2_M(\vth)+\mathrm{const}$.

\section{Подбор (структурный выбор) модели}
Даны конкурирующие структуры $\{M_1,\dots,M_m\}$ (физические гипотезы). Прямое сравнение по
$\chi^2$ некорректно: более сложная модель всегда подгоняет не хуже. Используем штраф за
сложность — информационный критерий Акаике:
\[
  \AIC(M)\;=\;2\,p_M\;-\;2\ln L_M(\hat{\vth}_M)
  \;=\;2\,p_M\;+\;\chi^2_M(\hat{\vth}_M)\;+\;\mathrm{const}.
\]
Выбор структуры и попарное сравнение:
\[
  \boxed{\;M^\star=\argmin_{M\in\mathcal{M}}\AIC(M)\;},
  \qquad
  \Delta\AIC(M)=\AIC(M)-\AIC(M^\star)\ \ge 0 .
\]
$\Delta\AIC\lesssim 2$ — модели неразличимы; $\Delta\AIC\gtrsim 10$ — уверенное отвержение.

\section{Идентифицируемость и вырожденность}
Даже при найденном $\hat{\vth}_M$ параметры могут быть не восстановимы по данным. Локальная
идентифицируемость определяется матрицей (Фишера) кривизны
\[
  \mathcal{I}(\vth)=J^{\mathsf H}WJ,\qquad
  J_{ka}=\frac{\partial \hat y_k}{\partial \vartheta_a},\qquad
  \mathrm{Cov}(\hat{\vth})\approx \mathcal{I}^{-1},
\]
а вырожденность — её плохой обусловленностью: если корреляция
$\rho_{ab}=\mathcal{I}^{-1}_{ab}/\sqrt{\mathcal{I}^{-1}_{aa}\mathcal{I}^{-1}_{bb}}\to\pm1$, пара
$(\vartheta_a,\vartheta_b)$ совместно неидентифицируема (в проекте это $\Deff\leftrightarrow\eta$,
$\rho\to\pm1$). Структурная идентифицируемость: $f(\cdot;\vth)\equiv f(\cdot;\vth')\Rightarrow
\vth=\vth'$.

\section{Оценка неопределённости (UQ)}
Каждому извлечённому параметру сопоставляется доверительный интервал. Профильное правдоподобие:
\[
  \mathrm{CI}_{1-\alpha}(\vartheta_a)=\Bigl\{\vartheta_a:\;
  \min_{\vth\setminus\vartheta_a}\chi^2_M(\vth)\;-\;\chi^2_{M,\min}\;\le\;\Delta_{1,\alpha}\Bigr\},
\]
где $\Delta_{1,\alpha}$ — квантиль $\chi^2_1$; параллельно — непараметрический бутстрап по
остаткам. Тесты вырождения обязательны там, где $\rho\to\pm1$.

\section{Целевой научный вопрос}
Внутри отобранной модели эффективный диаметр удовлетворяет $\Deff/\Dphys\approx 0.43\text{--}0.64$,
тогда как жёсткая привязка $\Deff=\Dphys$ катастрофична ($\Delta\AIC\gg 0$). Отсюда корректно
поставленная подзадача: \emph{идентифицировать физический канал, прокси которого служит} $\Deff$
(цель проекта №5).

\end{document}
